\documentclass[11pt]{book}
\usepackage[paperwidth=145mm,paperheight=200mm,top=18mm,bottom=20mm,inner=20mm,outer=16mm]{geometry}
\usepackage{brumfiel}
% figures01.tex -- Chapter 1 (INTRODUCTION) figures for Brumfiel, _Geometry_.
% Macro names: \FIG + I (chapter 1) + spelled-out figure number(s), letters only.
% Every constrained point is computed, never eyeballed:
%   bisectors  -> equal-angle construction / (intersection of ...)
%   altitudes  -> ($(A)!(P)!(B)$)
%   midpoints  -> ($(A)!0.5!(B)$)
%   parallels  -> shared direction vector / equal ratios
% Constraint audit: tools/constraints/ch01.py

%--------------------------------------------------------------- Fig 1-1
% Plane figures. Freeform illustrative shapes traced from scan01 p24 (book p.7).
\newcommand{\FIGIONE}{%
\begin{bkfigure}[\linewidth]
\begin{tikzpicture}[scale=0.92, every node/.style={font=\scriptsize}]
  % --- row 1
  \draw[given] (0,0) -- (2.0,0);
  \fill (0,0) circle (1.1pt) (2.0,0) circle (1.1pt);
  \node[vlab, above] at (0,0.02) {$A$};
  \node[vlab, above] at (2.0,0.02) {$B$};
  \node[below=5pt] at (1.0,-0.05) {Line segment};

  \draw[given] (3.3,0) -- (5.3,0);
  \fill (3.3,0) circle (1.1pt);
  \node[vlab, above] at (3.3,0.02) {$O$};
  \node[below=5pt] at (4.3,-0.05) {Ray from $O$};

  \draw[given] (6.6,-0.12) -- (8.7,-0.12);
  \draw[given] (6.6,-0.12) -- (8.25,0.95);
  \node[below=5pt] at (7.6,-0.17) {Angle};

  % --- row 2
  \begin{scope}[yshift=-2.35cm]
    \draw[given] (0.05,0.22) -- (0.95,1.05) -- (1.95,0) -- cycle;
    \node[below=5pt] at (1.0,-0.05) {Triangle};

    \draw[given] (2.85,0) rectangle (4.15,1.3);
    \node[below=5pt] at (3.5,-0.05) {Square};

    \draw[given] (4.95,0) -- (6.75,0) -- (7.35,1.15) -- (5.55,1.15) -- cycle;
    \node[below=5pt] at (6.15,-0.05) {Parallelogram};

    \draw[given] (8.6,0.6) circle (0.62);
    \node[below=5pt] at (8.6,-0.05) {Circle};
  \end{scope}

  % --- row 3
  \begin{scope}[yshift=-4.7cm]
    \draw[given] (0,0) -- (2.0,0) -- (1.65,1.15) -- (0.35,1.15) -- cycle;
    \node[below=5pt] at (1.0,-0.05) {Trapezoid};

    \draw[given] (3.05,0.62) -- (3.85,1.35) -- (5.05,0) -- (3.25,0.28) -- cycle;
    \node[align=center, below=5pt] at (4.0,-0.05)
         {Quadrilateral, or\\4-sided polygon};

    \draw[given] (6.5,0.55) -- (7.05,1.35) -- (8.35,0.62) -- (7.75,0) -- (6.95,0.1) -- cycle;
    \node[align=center, below=5pt] at (7.5,-0.05)
         {Pentagon, or\\5-sided polygon};
  \end{scope}

  % --- row 4
  \begin{scope}[yshift=-7.35cm]
    \draw[given] (3.3,0) rectangle (5.5,1.25);
    \node[below=5pt] at (4.4,-0.05) {Rectangle};
  \end{scope}
\end{tikzpicture}
\figcap{1--1}
\end{bkfigure}}

%--------------------------------------------------------------- Fig 1-2
% Solid figures. Hidden edges dashed, as the book draws them.
\newcommand{\FIGITWO}{%
\begin{bkfigure}[\linewidth]
\begin{tikzpicture}[scale=0.92, every node/.style={font=\scriptsize},
                    hid/.style={fig, dashed, dash pattern=on 2pt off 1.6pt}]
  % ---- cube
  \begin{scope}
    \def\d{0.42}
    \draw[given] (0,0) rectangle (1.25,1.25);
    \draw[given] (\d,\d) -- (\d+1.25,\d) -- (\d+1.25,\d+1.25) -- (\d,\d+1.25);
    \draw[hid]   (\d,\d) -- (\d,\d+1.25);
    \draw[given] (0,1.25) -- (\d,\d+1.25);
    \draw[given] (1.25,1.25) -- (\d+1.25,\d+1.25);
    \draw[given] (1.25,0) -- (\d+1.25,\d);
    \draw[hid]   (0,0) -- (\d,\d);
    \node[below=5pt] at (0.85,-0.05) {Cube};
  \end{scope}
  % ---- box
  \begin{scope}[xshift=3.15cm]
    \def\d{0.42}
    \draw[given] (0,0) rectangle (1.95,1.05);
    \draw[given] (\d,\d) -- (\d+1.95,\d) -- (\d+1.95,\d+1.05) -- (\d,\d+1.05);
    \draw[hid]   (\d,\d) -- (\d,\d+1.05);
    \draw[given] (0,1.05) -- (\d,\d+1.05);
    \draw[given] (1.95,1.05) -- (\d+1.95,\d+1.05);
    \draw[given] (1.95,0) -- (\d+1.95,\d);
    \draw[hid]   (0,0) -- (\d,\d);
    \node[align=center, below=5pt] at (1.2,-0.05)
         {Rectangular\\parallelepiped or box};
  \end{scope}
  % ---- circular cylinder
  \begin{scope}[xshift=6.95cm]
    \draw[given] (0,0) -- (0,1.6);
    \draw[given] (1.3,0) -- (1.3,1.6);
    \draw[given] (0.65,1.6) ellipse (0.65 and 0.24);
    \draw[given] (0.65,0) ++(180:0.65) arc (180:360:0.65 and 0.24);
    \draw[hid]   (0.65,0) ++(180:0.65) arc (180:0:0.65 and 0.24);
    \node[align=center, below=5pt] at (0.65,-0.26) {Circular\\cylinder};
  \end{scope}

  % ---- tetrahedron: as the book draws it, the near vertex is at the front and
  % the single hidden element is the back base edge L--R
  \begin{scope}[yshift=-3.1cm]
    \coordinate (tL) at (0,0.402);
    \coordinate (tR) at (1.517,0.674);
    \coordinate (tF) at (0.817,0);
    \coordinate (tT) at (0.596,1.750);
    \draw[given] (tT) -- (tL) -- (tF) -- (tR) -- cycle;
    \draw[given] (tT) -- (tF);
    \draw[hid]   (tL) -- (tR);
    \node[below=5pt] at (0.76,-0.1) {Tetrahedron};
  \end{scope}
  % ---- pyramid
  \begin{scope}[xshift=3.15cm, yshift=-3.1cm]
    \coordinate (p1) at (0,0.15);
    \coordinate (p2) at (1.45,0);
    \coordinate (p3) at (2.05,0.55);
    \coordinate (p4) at (0.6,0.72);
    \coordinate (pa) at (1.05,1.75);
    \draw[given] (p1) -- (p2) -- (p3);
    \draw[hid]   (p3) -- (p4) -- (p1);
    \draw[given] (p1) -- (pa) -- (p2);
    \draw[given] (p3) -- (pa);
    \draw[hid]   (p4) -- (pa);
    \node[below=5pt] at (1.0,-0.1) {Pyramid};
  \end{scope}
  % ---- prism: upright triangular prism, near vertical edge facing the reader,
  % top face fully visible, only the back bottom edge hidden (as in the book)
  \begin{scope}[xshift=6.95cm, yshift=-3.1cm]
    \coordinate (uL) at (0,1.732);        \coordinate (dL) at (0,0.426);
    \coordinate (uR) at (1.205,1.656);    \coordinate (dR) at (1.205,0.350);
    \coordinate (uF) at (0.266,1.306);    \coordinate (dF) at (0.266,0);
    \draw[given] (uL) -- (uR) -- (uF) -- cycle;   % top face
    \draw[given] (uL) -- (dL);
    \draw[given] (uR) -- (dR);
    \draw[given] (uF) -- (dF);
    \draw[given] (dL) -- (dF) -- (dR);
    \draw[hid]   (dL) -- (dR);
    \node[below=5pt] at (0.60,-0.1) {Prism};
  \end{scope}

  % ---- sphere
  \begin{scope}[yshift=-6.2cm]
    \draw[given] (0.85,0.85) circle (0.78);
    \draw[given] (0.85,0.85) ++(180:0.78) arc (180:360:0.78 and 0.27);
    \draw[hid]   (0.85,0.85) ++(180:0.78) arc (180:0:0.78 and 0.27);
    \node[below=5pt] at (0.85,-0.05) {Sphere};
  \end{scope}
  % ---- torus
  \begin{scope}[xshift=3.15cm, yshift=-6.2cm]
    \draw[given] (1.05,0.85) circle (0.82);
    \draw[given] (1.12,0.92) circle (0.40);
    \draw[given] (0.72,0.92) arc (180:300:0.40 and 0.16);
    \node[align=center, below=5pt] at (1.05,-0.05)
         {Torus, or doughnut,\\or anchor ring};
  \end{scope}
  % ---- cone
  \begin{scope}[xshift=6.95cm, yshift=-6.2cm]
    \draw[given] (0,0.35) -- (0.78,1.85) -- (1.56,0.35);
    \draw[given] (0.78,0.35) ++(180:0.78) arc (180:360:0.78 and 0.27);
    \draw[hid]   (0.78,0.35) ++(180:0.78) arc (180:0:0.78 and 0.27);
    \node[below=5pt] at (0.78,-0.05) {Cone};
  \end{scope}
\end{tikzpicture}
\figcap{1--2}
\end{bkfigure}}

%--------------------------------------------------------------- Fig 1-3
% "Strange" sets of points. ORGANIC -- traced against scan02 p2 (book p.9):
% damped zigzag, chain of shrinking circles, recursively quartered square,
% face built from circles.
\newcommand{\FIGITHREE}{%
\begin{bkfigure}[0.92\linewidth]
\begin{tikzpicture}[scale=1.0]
  % ---- damped zigzag: amplitude falls geometrically left to right
  \begin{scope}[shift={(0,3.05)}]
    % the path must OPEN with a coordinate -- starting it with "--" emits a
    % lineto with no current point and some viewers drop the whole zigzag
    \draw[given] (0,0.42)
      \foreach \i [evaluate={\x=0.30*\i; \s=ifthenelse(mod(\i,2)==0,1,-1);
                            \a=0.42*pow(0.80,\i)} ] in {1,...,15} { -- (\x,{\s*\a}) }
      -- (4.85,0);
  \end{scope}
  % ---- chain of shrinking circles converging to a point
  \begin{scope}[shift={(0.15,1.55)}]
    \foreach \i [evaluate={\r=0.235*pow(0.80,\i); \c=0.44*(1-pow(0.80,\i))/0.20}] in {0,...,9}
      \draw[given] ({\c+\r},0) ellipse ({\r} and {1.28*\r});
    \draw[given] (2.30,0.05) -- (2.62,0) -- (2.30,-0.05);
  \end{scope}
  % ---- square quartered again and again in its upper-right corner
  \begin{scope}[shift={(3.15,0.60)}]
    \def\S{2.0}
    \draw[given] (0,0) rectangle (\S,\S);
    \draw[given] (0,{\S/2}) -- (\S,{\S/2});
    \draw[given] ({\S/2},0) -- ({\S/2},\S);
    \draw[given] ({\S/2},{3*\S/4}) -- (\S,{3*\S/4});
    \draw[given] ({3*\S/4},{\S/2}) -- ({3*\S/4},\S);
    \draw[given] ({3*\S/4},{7*\S/8}) -- (\S,{7*\S/8});
    \draw[given] ({7*\S/8},{3*\S/4}) -- ({7*\S/8},\S);
    \draw[given] ({7*\S/8},{15*\S/16}) -- (\S,{15*\S/16});
    \draw[given] ({15*\S/16},{7*\S/8}) -- ({15*\S/16},\S);
  \end{scope}
  % ---- face made of circles
  \begin{scope}[shift={(1.55,-1.25)}]
    \draw[given] (0,0) circle (0.80 and 0.86);      % head
    \draw[given] (-0.86,0.02) circle (0.10);        % ears
    \draw[given] ( 0.86,0.02) circle (0.10);
    \draw[given] (-0.30,0.26) circle (0.15);        % eyes
    \draw[given] ( 0.30,0.26) circle (0.15);
    \draw[given] (-0.30,0.26) circle (0.06);
    \draw[given] ( 0.30,0.26) circle (0.06);
    \draw[given] (0,-0.06) circle (0.09 and 0.14);  % nose
    \draw[given] (0,-0.48) circle (0.17);           % mouth
  \end{scope}
\end{tikzpicture}
\figcap{1--3}
\end{bkfigure}}

%--------------------------------------------------------------- Fig 1-4
% Angle: vertex, interior, exterior.
\newcommand{\FIGIFOUR}{%
\begin{bkfigure}[0.60\linewidth]
\begin{tikzpicture}[scale=1.0]
  \coordinate (A) at (0,0);
  \draw[given] (A) -- (2.35,1.20);
  \draw[given] (A) -- (2.35,-1.20);
  \node[vlab] at (-0.20,0) {$A$};
  \node[slab] at (-0.42,0.42) {Vertex};
  \node[slab] at (1.02,0) {Interior};
  \node[slab] at (0.62,-0.88) {Exterior};
\end{tikzpicture}
\figcap{1--4}
\end{bkfigure}}

%--------------------------------------------------------------- Figs 1-5, 1-6
% 1-5: ray AC bisects angle BAD -- C placed at exactly half the angle of B.
% 1-6: acute (< 90) and obtuse (> 90) angles.
\newcommand{\FIGIFIVESIX}{%
\begin{bkfigure}[\linewidth]
\begin{minipage}[t]{0.48\linewidth}\centering
\begin{tikzpicture}[scale=1.0]
  \coordinate (A) at (0,0);
  \coordinate (D) at (2.30,0);
  \coordinate (B) at (58:2.30);
  \coordinate (C) at (29:2.15);       % exact bisector of angle BAD
  \draw[given] (A) -- (D);
  \draw[given] (A) -- (B);
  \draw[given] (A) -- (C);
  \node[vlab, left]  at (A) {$A$};
  \node[vlab, right] at (D) {$D$};
  \node[vlab, above] at (B) {$B$};
  \node[vlab, right] at (C) {$C$};
\end{tikzpicture}
\figcap{1--5}
\end{minipage}\hfill
\begin{minipage}[t]{0.48\linewidth}\centering
\begin{tikzpicture}[scale=1.0]
  % acute: vertex bottom-left, 32 degrees
  \draw[given] (0,0) -- (1.55,0);
  \draw[given] (0,0) -- (32:1.55);
  \node[slab, above] at (0.72,0.85) {Acute angle};
  % obtuse: vertex bottom-right; leftward ray at 180 deg, second ray at 47 deg,
  % so the angle at the vertex is 180 - 47 = 133 deg
  \draw[given] (3.35,0) -- (1.95,0);
  \draw[given] (3.35,0) -- ($(3.35,0)+(47:1.55)$);
  \node[slab, above] at (2.85,1.15) {Obtuse angle};
\end{tikzpicture}
\figcap{1--6}
\end{minipage}
\end{bkfigure}}

%--------------------------------------------------------------- Figs 1-7, 1-8
% 1-7: right triangle, right angle at the lower-left vertex.
% 1-8: angle AOB is a right angle; C interior.
\newcommand{\FIGISEVENEIGHT}{%
\begin{bkfigure}[\linewidth]
\begin{minipage}[t]{0.48\linewidth}\centering
\begin{tikzpicture}[scale=1.0]
  \coordinate (R) at (0,0);
  \coordinate (T) at (0,1.75);
  \coordinate (S) at (2.15,0);
  \draw[given] (R) -- (T) -- (S) -- cycle;
  \draw[fig] (0,0.17) -- (0.17,0.17) -- (0.17,0);
  \node[slab] at (0.56,0.30) {90\textdegree};
\end{tikzpicture}
\figcap{1--7}
\end{minipage}\hfill
\begin{minipage}[t]{0.48\linewidth}\centering
\begin{tikzpicture}[scale=1.0]
  \coordinate (O) at (0,0);
  \coordinate (A) at (0,1.75);
  \coordinate (B) at (1.85,0);
  \coordinate (C) at (33:1.62);
  \draw[given] (O) -- (A);
  \draw[given] (O) -- (B);
  \draw[given] (O) -- (C);
  \fill (A) circle (1.2pt) (B) circle (1.2pt) (C) circle (1.2pt);
  \node[vlab, above] at (A) {$A$};
  \node[vlab, right] at (B) {$B$};
  \node[vlab, right] at (C) {$C$};
  \node[vlab, below left, xshift=-1pt, yshift=-0pt] at (O) {$O$};
\end{tikzpicture}
\figcap{1--8}
\end{minipage}
\end{bkfigure}}

%--------------------------------------------------------------- Figs 1-9, 1-10
% 1-9: A, O, C collinear on l (so angle AOB and angle BOC are supplementary).
% 1-10: straight angle A-O-B.
\newcommand{\FIGININETEN}{%
\begin{bkfigure}[\linewidth]
\begin{minipage}[t]{0.48\linewidth}\centering
\begin{tikzpicture}[scale=1.0]
  \coordinate (A) at (-1.55,0);
  \coordinate (O) at (0,0);
  \coordinate (C) at (1.55,0);
  \coordinate (B) at (38:1.62);
  \draw[given] (A) -- (C);
  \draw[given] (O) -- (B);
  \fill (A) circle (1.2pt) (O) circle (1.2pt) (C) circle (1.2pt) (B) circle (1.2pt);
  \node[vlab, left]  at (A) {$A$};
  \node[vlab, right] at (C) {$C$};
  \node[vlab, above right] at (B) {$B$};
  \node[vlab, below] at (O) {$O$};
  \node[vlab, below] at (0.80,-0.02) {$l$};
\end{tikzpicture}
\figcap{1--9}
\end{minipage}\hfill
\begin{minipage}[t]{0.48\linewidth}\centering
\begin{tikzpicture}[scale=1.0]
  \coordinate (A) at (-1.45,0);
  \coordinate (O) at (0,0);
  \coordinate (B) at (1.45,0);
  \draw[given] (A) -- (B);
  \fill (A) circle (1.2pt) (O) circle (1.2pt) (B) circle (1.2pt);
  \node[vlab, above] at (A) {$A$};
  \node[vlab, above] at (B) {$B$};
  \node[vlab, below] at (O) {$O$};
\end{tikzpicture}
\figcap{1--10}
\end{minipage}
\end{bkfigure}}

%--------------------------------------------------------------- Fig 1-11
\newcommand{\FIGIELEVEN}{%
\begin{bkfigure}[0.52\linewidth]
\begin{tikzpicture}[scale=1.0]
  % vertex ratios measured off the book's own quadrilateral
  \draw[given] (0,0.310) -- (0.558,1.565) -- (2.200,0.558) -- (1.270,0) -- cycle;
\end{tikzpicture}
\figcap{1--11}
\end{bkfigure}}

%--------------------------------------------------------------- Figs 1-12, 1-13
% 1-12: equilateral triangle by the two-arc construction (Euclid I.1).
% 1-13: an angle and its copy -- both drawn at exactly 52 degrees.
\newcommand{\FIGITWELVETHIRTEEN}{%
\begin{bkfigure}[\linewidth]
\begin{minipage}[t]{0.46\linewidth}\centering
\begin{tikzpicture}[scale=1.0]
  \def\s{2.05}
  \coordinate (A) at (0,0);
  \coordinate (B) at (\s,0);
  \coordinate (C) at (60:\s);          % exact equilateral apex
  \draw[given] (A) -- (B) -- (C) -- cycle;
  \draw[fig] (A) ++(-18:\s) arc (-18:78:\s);
  \draw[fig] (B) ++(198:\s) arc (198:102:\s);
  \draw[fig] ($(C)+(-0.09,-0.09)$) -- ($(C)+(0.09,0.09)$);
  \draw[fig] ($(C)+(-0.09,0.09)$) -- ($(C)+(0.09,-0.09)$);
  % compass-point marks where each arc passes through the opposite base vertex
  \draw[fig] ($(A)+(135:0.11)$) -- ($(A)+(-45:0.11)$);
  \draw[fig] ($(B)+(45:0.11)$)  -- ($(B)+(225:0.11)$);
\end{tikzpicture}
\figcap{1--12}
\end{minipage}\hfill
\begin{minipage}[t]{0.50\linewidth}\centering
\begin{tikzpicture}[scale=1.0]
  \foreach \x/\lo/\lb/\la in {0/{O}/{B}/{A}, 2.75/{O\mkern2mu'}/{B\mkern2mu'}/{A\mkern2mu'}} {
    \begin{scope}[xshift=\x cm]
      \coordinate (V) at (0,0);
      \draw[given] (V) -- (1.45,0);
      \draw[given] (V) -- (60:1.60);
      \draw[fig] (V) ++(0:0.62) arc (0:60:0.62);
      % the book marks where the arc cuts each ray: a tick on OB, a cross on OA
      \draw[fig] ($(V)+(0.62,-0.10)$) -- ($(V)+(0.62,0.10)$);
      \draw[fig] ($(V)+(60:0.62)+(15:0.11)$)  -- ($(V)+(60:0.62)+(195:0.11)$);
      \draw[fig] ($(V)+(60:0.62)+(105:0.11)$) -- ($(V)+(60:0.62)+(285:0.11)$);
      \node[vlab, below left, xshift=-0pt, yshift=-0pt] at (V) {$\lo$};
      \node[vlab, right] at (1.45,0) {$\lb$};
      \node[vlab, above] at (60:1.60) {$\la$};
    \end{scope}}
\end{tikzpicture}
\figcap{1--13}
\end{minipage}
\end{bkfigure}}

%--------------------------------------------------------------- Fig 1-14
% Compass designs. ORGANIC -- traced against scan02 p9 (book p.14):
% (left) four circles of radius r whose centres lie at distance r from a common
% point; (right) the six-petal rosette inscribed in a circle.
\newcommand{\FIGIFOURTEEN}{%
\begin{bkfigure}[0.86\linewidth]
\begin{tikzpicture}[scale=1.0]
  % four-circle design
  \begin{scope}
    \def\r{0.78}
    \foreach \a in {0,90,180,270} \draw[given] (\a:\r) circle (\r);
  \end{scope}
  % six-petal rosette
  \begin{scope}[xshift=4.15cm]
    \def\R{1.12}
    \begin{scope}
      \clip (0,0) circle (\R);
      \foreach \a in {0,60,...,300} \draw[given] (\a:\R) circle (\R);
    \end{scope}
    \draw[given] (0,0) circle (\R);
  \end{scope}
\end{tikzpicture}
\figcap{1--14}
\end{bkfigure}}

%--------------------------------------------------------------- Figs 1-15, 1-16
% 1-15: angle bisected with compass; Z is the exact arc crossing, so VZ bisects.
% 1-16: adjacent angles AOB and BOC sharing ray OB.
\newcommand{\FIGIFIFTEENSIXTEEN}{%
\begin{bkfigure}[\linewidth]
\begin{minipage}[t]{0.48\linewidth}\centering
\begin{tikzpicture}[scale=1.0]
  \coordinate (V) at (0,0);
  \draw[given] (V) -- (26:2.15);
  \draw[given] (V) -- (-26:2.15);
  \coordinate (X) at (26:0.72);
  \coordinate (Y) at (-26:0.72);
  \draw[fig] (V) ++(34:0.72) arc (34:-34:0.72);
  % arcs of radius 0.716170 about X and Y meet at Z on the bisector; as in the
  % book only their crossing is marked, the two small arcs are not drawn
  \coordinate (Z) at (1.29,0);
  \draw[given] (V) -- (2.15,0);
  \draw[fig] ($(Z)+(-0.09,-0.09)$) -- ($(Z)+(0.09,0.09)$);
  \draw[fig] ($(Z)+(-0.09,0.09)$) -- ($(Z)+(0.09,-0.09)$);
\end{tikzpicture}
\figcap{1--15}
\end{minipage}\hfill
\begin{minipage}[t]{0.48\linewidth}\centering
\begin{tikzpicture}[scale=1.0]
  \coordinate (O) at (0,0);
  \coordinate (A) at (36:2.05);
  \coordinate (B) at (-2:2.05);
  \coordinate (C) at (-24:2.05);
  \draw[given] (O) -- (A);
  \draw[given] (O) -- (B);
  \draw[given] (O) -- (C);
  \node[vlab, left]  at (O) {$O$};
  \node[vlab, above right] at (A) {$A$};
  \node[vlab, right] at (B) {$B$};
  \node[vlab, right] at (C) {$C$};
\end{tikzpicture}
\figcap{1--16}
\end{minipage}
\end{bkfigure}}

%--------------------------------------------------------------- Figs 1-17, 1-18
% 1-17: three pairs of angles, none of them adjacent (traced from scan02 p9):
%   (a) two angles at two DIFFERENT vertices;
%   (b) two angles sharing a vertex and a side but overlapping interiors;
%   (c) two angles sharing only the vertex -- no common side.
% 1-18: copying a segment onto a line at A.
\newcommand{\FIGISEVENTEENEIGHTEEN}{%
\begin{bkfigure}[\linewidth]
\begin{minipage}[t]{0.56\linewidth}\centering
\begin{tikzpicture}[scale=0.95]
  % (a) two nearby but distinct vertices
  \begin{scope}
    \coordinate (V) at (0,0);
    \coordinate (W) at (0.16,-0.14);
    \draw[given] (V) -- (34:1.75);
    \draw[given] (V) -- (-9:1.70);
    \draw[given] (W) -- ($(W)+(-6:1.55)$);
    \draw[given] (W) -- ($(W)+(-44:1.30)$);
    \node[vlab] at (12.5:0.80) {$A$};
    \node[vlab] at ($(W)+(-25:0.98)$) {$B$};
  \end{scope}
  % (b) shared vertex and side, interiors overlap
  \begin{scope}[xshift=3.35cm, yshift=0.05cm]
    \coordinate (V) at (1.30,0);
    \draw[given] (V) -- ($(V)+(140:1.45)$);
    \draw[given] (V) -- ($(V)+(218:1.35)$);
    \draw[given] (V) -- ($(V)+(2:1.15)$);
    \node[vlab] at ($(V)+(35:0.52)$) {$A$};
    \node[vlab] at ($(V)+(-35:0.52)$) {$B$};
  \end{scope}
  % (c) shared vertex only
  \begin{scope}[xshift=1.55cm, yshift=-2.55cm]
    \coordinate (V) at (0,0);
    \draw[given] (V) -- (92:1.55);
    \draw[given] (V) -- (62:1.45);
    \draw[given] (V) -- (-10:1.50);
    \draw[given] (V) -- (-38:1.45);
    \node[vlab] at (77:1.15) {$A$};
    \node[vlab] at (-24:1.25) {$B$};
  \end{scope}
\end{tikzpicture}
\figcap{1--17}
\end{minipage}\hfill
\begin{minipage}[t]{0.40\linewidth}\centering
\begin{tikzpicture}[scale=1.0]
  \draw[given] (0.28,0.52) -- (2.05,0.52);        % the given segment
  \draw[given] (0,0) -- (2.30,0);                  % the line
  \coordinate (A) at (0.28,0);
  \fill (A) circle (1.2pt);
  \draw[fig] (2.05,-0.13) -- (2.05,0.13);          % compass mark at B
  \node[vlab, below left, xshift=2pt, yshift=-0pt] at (A) {$A$};
  \node[vlab, above right, xshift=-3pt, yshift=0pt] at (2.05,0.02) {$B$};
\end{tikzpicture}
\figcap{1--18}
\end{minipage}
\end{bkfigure}}

%--------------------------------------------------------------- Figs 1-19, 1-20
% 1-19: perpendicular bisector of AB; C is the exact midpoint, l is exactly
%       perpendicular (vertical through C).
% 1-20: perpendicular to l at a point O on l.
\newcommand{\FIGININETEENTWENTY}{%
\begin{bkfigure}[\linewidth]
\begin{minipage}[t]{0.48\linewidth}\centering
\begin{tikzpicture}[scale=1.0]
  % drawn 1.4x the first draft's size, to the book's own proportions
  \coordinate (A) at (-1.89,0);
  \coordinate (B) at (1.89,0);
  \coordinate (C) at ($(A)!0.5!(B)$);     % exact midpoint
  \draw[given] (A) -- (B);
  \draw[given] (0,-2.17) -- (0,2.17);
  % compass arcs of radius 2.45 about A and B; the sweep runs past each
  % crossing, as the book draws it. Chained 23-degree pieces keep every
  % Bezier tight against its own arc.
  \draw[aux] (A) ++(-46:2.45) arc (-46:-23:2.45) arc (-23:0:2.45)
                              arc (0:23:2.45)    arc (23:46:2.45);
  \draw[aux] (B) ++(134:2.45) arc (134:157:2.45) arc (157:180:2.45)
                              arc (180:203:2.45) arc (203:226:2.45);
  \foreach \y in {1.558974,-1.558974} {   % sqrt(2.45^2 - 1.89^2)
    \draw[fig] (-0.12,\y-0.12) -- (0.12,\y+0.12);
    \draw[fig] (-0.12,\y+0.12) -- (0.12,\y-0.12);}
  % where the arcs cut AB (x = -1.89+2.45 and 1.89-2.45)
  \foreach \x in {-0.56,0.56} \draw[fig] (\x,-0.15) -- (\x,0.15);
  \fill (A) circle (1.2pt) (B) circle (1.2pt);
  \node[vlab, left]  at (A) {$A$};
  \node[vlab, right] at (B) {$B$};
  \node[vlab] at (0.20,-0.33) {$C$};
  \node[vlab, below right, xshift=0pt, yshift=-1pt] at (0,-2.17) {$l$};
\end{tikzpicture}
\figcap{1--19}
\end{minipage}\hfill
\begin{minipage}[t]{0.48\linewidth}\centering
\begin{tikzpicture}[scale=1.0]
  % same 1.4x enlargement as Fig 1-19, so the pair matches
  \coordinate (O) at (0,0);
  \draw[given] (-2.03,0) -- (2.03,0);
  \draw[given] (0,-2.03) -- (0,2.17);
  \foreach \x in {-1.12,1.12} \draw[fig] (\x,-0.18) -- (\x,0.18);
  \draw[aux] (-1.12,0) ++(-58:1.61)
      arc (-58:-38.667:1.61) arc (-38.667:-19.333:1.61) arc (-19.333:0:1.61)
      arc (0:19.333:1.61)    arc (19.333:38.667:1.61)   arc (38.667:58:1.61);
  \draw[aux] ( 1.12,0) ++(122:1.61)
      arc (122:141.333:1.61)   arc (141.333:160.667:1.61) arc (160.667:180:1.61)
      arc (180:199.333:1.61)   arc (199.333:218.667:1.61) arc (218.667:238:1.61);
  \foreach \y in {1.156590,-1.156590} {   % sqrt(1.61^2 - 1.12^2)
    \draw[fig] (-0.12,\y-0.12) -- (0.12,\y+0.12);
    \draw[fig] (-0.12,\y+0.12) -- (0.12,\y-0.12);}
  \node[vlab, above right, xshift=0pt, yshift=0pt] at (0,2.17) {$m$};
  \node[vlab, left] at (-2.03,0) {$l$};
  \node[vlab] at (0.20,-0.33) {$O$};
\end{tikzpicture}
\figcap{1--20}
\end{minipage}
\end{bkfigure}}

%--------------------------------------------------------------- Figs 1-21, 1-22
% 1-21: perpendicular to l from an external point P (m is exactly vertical).
% 1-22: the three medians; O is the exact centroid.
\newcommand{\FIGITWENTYONETWENTYTWO}{%
\begin{bkfigure}[\linewidth]
\begin{minipage}[t]{0.46\linewidth}\centering
\begin{tikzpicture}[scale=1.0]
  \coordinate (P) at (0,1.30);
  \draw[given] (-1.60,0) -- (1.60,0);
  \draw[given] (0,-1.15) -- (0,1.55);
  % arc of radius 1.75 about P cuts l at x = +- sqrt(1.75^2 - 1.30^2) = 1.171538;
  % the crossings are at 227.95 deg and 312.05 deg, so the sweep runs past them
  \draw[fig] (P) ++(222:1.75) arc (222:318:1.75);
  \draw[fig] (-0.09,1.30-0.09) -- (0.09,1.30+0.09);
  \draw[fig] (-0.09,1.30+0.09) -- (0.09,1.30-0.09);
  % arcs of radius 1.45 about those points meet at y = -sqrt(1.45^2 - 1.171538^2)
  \draw[fig] (-0.09,-0.854400-0.09) -- (0.09,-0.854400+0.09);
  \draw[fig] (-0.09,-0.854400+0.09) -- (0.09,-0.854400-0.09);
  \node[vlab, left] at (P) {$P$};
  \node[vlab, left] at (-1.60,0) {$l$};
  \node[vlab, below] at (0,-1.15) {$m$};
\end{tikzpicture}
\figcap{1--21}
\end{minipage}\hfill
\begin{minipage}[t]{0.50\linewidth}\centering
\begin{tikzpicture}[scale=1.0]
  % proportions from the book: apex at 0.64 of the base, height 0.627 of the base
  \coordinate (A) at (0,0);
  \coordinate (B) at (3.40,0);
  \coordinate (C) at (2.176,2.132);
  \coordinate (Mc) at ($(A)!0.5!(B)$);
  \coordinate (Ma) at ($(B)!0.5!(C)$);
  \coordinate (Mb) at ($(C)!0.5!(A)$);
  \coordinate (O)  at ($(A)!0.6666666!(Ma)$);   % exact centroid
  \draw[given] (A) -- (B) -- (C) -- cycle;
  \draw[aux] (A) -- (Ma);
  \draw[aux] (B) -- (Mb);
  \draw[aux] (C) -- (Mc);
  \draw[fig] ($(O)+(-0.08,-0.08)$) -- ($(O)+(0.08,0.08)$);
  \draw[fig] ($(O)+(-0.08,0.08)$) -- ($(O)+(0.08,-0.08)$);
  % 296.33 deg bisects the widest gap between the six median rays at O
  \node[vlab] at ($(O)+(296.33:0.47)$) {$O$};
\end{tikzpicture}
\figcap{1--22}
\end{minipage}
\end{bkfigure}}

%--------------------------------------------------------------- Fig 1-23
% Cardboard triangle balancing on a fingertip. ORGANIC (pictorial) -- traced
% against scan02 p12 (book p.16). The medians are exact; the hand is stylised.
\newcommand{\FIGITWENTYTHREE}{%
\begin{bkfigure}[0.50\linewidth]
\begin{tikzpicture}[scale=1.0]
  % triangle seen almost edge-on, balanced at its centroid
  \coordinate (A) at (-1.55,0.18);
  \coordinate (B) at (1.62,0.30);
  \coordinate (C) at (-0.28,0.92);
  \coordinate (Ma) at ($(B)!0.5!(C)$);
  \coordinate (Mb) at ($(C)!0.5!(A)$);
  \coordinate (Mc) at ($(A)!0.5!(B)$);
  \coordinate (O)  at ($(A)!0.6666666!(Ma)$);
  \draw[given] (A) -- (B) -- (C) -- cycle;
  \draw[aux] (A) -- (Ma);
  \draw[aux] (B) -- (Mb);
  \draw[aux] (C) -- (Mc);
  % stylised finger and hand: the rounded fingertip touches the centroid and
  % the finger runs down to meet the top of the fist (it must not float free)
  \draw[fig] ($(O)+(-0.075,-1.14)$) -- ($(O)+(-0.075,-0.02)$)
             arc (180:0:0.075) -- ($(O)+(0.075,-1.14)$);
  \draw[fig] (-0.30,-0.62) .. controls (-0.34,-0.98) and (-0.30,-1.24) .. (-0.24,-1.42);
  \draw[fig] ( 0.16,-0.62) .. controls (0.22,-0.98) and (0.24,-1.24) .. (0.20,-1.42);
  \draw[fig] (-0.30,-0.62) .. controls (-0.20,-0.70) and (0.06,-0.70) .. (0.16,-0.62);
  \foreach \k in {0,1,2} {
    \draw[fig] (-0.30,-0.78-0.17*\k)
       .. controls (-0.40,-0.84-0.17*\k) and (-0.42,-0.92-0.17*\k) ..
          (-0.32,-0.96-0.17*\k);}
  \draw[fig] (-0.30,-1.42) rectangle (0.22,-1.60);
\end{tikzpicture}
\figcap{1--23}
\end{bkfigure}}

%--------------------------------------------------------------- Fig 1-24
% Altitudes. (a) acute, foot inside; (b) obtuse, foot on the extension of the
% base; (c) the three altitude lines concurrent at O. All feet are exact
% projections ($(X)!(P)!(Y)$).
\newcommand{\FIGITWENTYFOUR}{%
\begin{bkfigure}[\linewidth]
\begin{tikzpicture}[scale=1.0]
  % Proportions measured off the book: all three panels are about the same width
  % (100 : 106 : 110) and grow in height (64 : 92 : 113).
  % ---- (a) apex left of centre, foot inside, right-angle box right of the altitude
  \begin{scope}
    \coordinate (A) at (0,0);
    \coordinate (B) at (2.15,0);
    \coordinate (C) at (0.7256,1.376);
    \coordinate (F) at ($(A)!(C)!(B)$);
    \draw[given] (A) -- (B) -- (C) -- cycle;
    \draw[aux] (C) -- (F);
    \draw[fig] ($(F)+(0.16,0)$) -- ($(F)+(0.16,0.16)$) -- ($(F)+(0,0.16)$);
    \node[slab, below=6pt] at (1.075,0) {(a)};
  \end{scope}
  % ---- (b) obtuse at B: the foot falls beyond B
  \begin{scope}[xshift=3.05cm]
    \coordinate (A) at (0,0);
    \coordinate (B) at (1.3975,0);
    \coordinate (C) at (2.2844,1.9834);
    \coordinate (F) at ($(A)!(C)!(B)$);
    \draw[given] (A) -- (B) -- (C) -- cycle;
    \draw[given] (B) -- (F);
    \draw[aux] (C) -- (F);
    \draw[fig] ($(F)+(-0.16,0)$) -- ($(F)+(-0.16,0.16)$) -- ($(F)+(0,0.16)$);
    \node[slab, below=6pt] at (1.142,0) {(b)};
  \end{scope}
  % ---- (c) three altitudes concurrent; near-isosceles, apex at 0.584 of the base
  \begin{scope}[xshift=6.10cm]
    \coordinate (A) at (0,0);
    \coordinate (B) at (2.365,0);
    \coordinate (C) at (1.3812,2.4295);
    \coordinate (Fc) at ($(A)!(C)!(B)$);
    \coordinate (Fa) at ($(B)!(A)!(C)$);
    \coordinate (Fb) at ($(A)!(B)!(C)$);
    \coordinate (H)  at (intersection of C--Fc and A--Fa);
    \draw[given] (A) -- (B) -- (C) -- cycle;
    \draw[aux] (C) -- (Fc);
    \draw[aux] (A) -- (Fa);
    \draw[aux] (B) -- (Fb);
    \draw[fig] ($(H)+(-0.08,-0.08)$) -- ($(H)+(0.08,0.08)$);
    \draw[fig] ($(H)+(-0.08,0.08)$) -- ($(H)+(0.08,-0.08)$);
    % 176.21 deg bisects the gap between the altitude rays at 150.38 and 202.05
    \node[vlab] at ($(H)+(176.21:0.65)$) {$O$};
    \node[slab, below=6pt] at (1.1825,0) {(c)};
  \end{scope}
\end{tikzpicture}
\figcap{1--24}
\end{bkfigure}}

%--------------------------------------------------------------- Fig 1-25
% The three angle bisectors meet at the incentre. Each bisector is drawn from a
% vertex through the exact incentre to the opposite side.
\newcommand{\FIGITWENTYFIVE}{%
\begin{bkfigure}[0.52\linewidth]
\begin{tikzpicture}[scale=1.0]
  % proportions from the book: apex at 0.775 of the base, height 0.667 of the base
  \coordinate (A) at (0,0);
  \coordinate (B) at (2.95,0);
  \coordinate (C) at (2.28625,1.96765);
  % incentre = (a A + b B + c C)/(a+b+c), a=|BC|, b=|CA|, c=|AB|
  % |BC| = 2.077003, |CA| = 3.016421, |AB| = 2.95  ->  I = (1.944900, 0.721694)
  \coordinate (I) at (1.944900,0.721694);
  \draw[given] (A) -- (B) -- (C) -- cycle;
  \draw[aux] (A) -- (intersection of A--I and B--C);
  \draw[aux] (B) -- (intersection of B--I and C--A);
  \draw[aux] (C) -- (intersection of C--I and A--B);
  \draw[fig] ($(I)+(-0.08,-0.08)$) -- ($(I)+(0.08,0.08)$);
  \draw[fig] ($(I)+(-0.08,0.08)$) -- ($(I)+(0.08,-0.08)$);
  % 227.52 deg bisects the gap between the bisector rays at 200.36 and 254.68
  \node[vlab] at ($(I)+(227.52:0.58)$) {$O$};
\end{tikzpicture}
\figcap{1--25}
\end{bkfigure}}

%--------------------------------------------------------------- Figs 1-26, 1-27
% 1-26: l parallel to m (same direction vector), transversal n; angles 1 and 2
%       are corresponding angles, hence equal.
% 1-27: copying angle 2 at P produces a line through P parallel to l.
\newcommand{\FIGITWENTYSIXTWENTYSEVEN}{%
\begin{bkfigure}[\linewidth]
\begin{minipage}[t]{0.48\linewidth}\centering
\begin{tikzpicture}[scale=1.0]
  % l and m share the direction (3.2, 0.9); n is the transversal
  \coordinate (l1) at (0.10,0.95);  \coordinate (l2) at (3.30,1.85);
  \coordinate (m1) at (0.10,0.30);  \coordinate (m2) at (3.30,1.20);
  \coordinate (n1) at (0.55,2.05);  \coordinate (n2) at (2.05,-0.10);
  \draw[given] (l1) -- (l2);
  \draw[given] (m1) -- (m2);
  \draw[given] (n1) -- (n2);
  \coordinate (P1) at (intersection of l1--l2 and n1--n2);
  \coordinate (P2) at (intersection of m1--m2 and n1--n2);
  % the book marks the two ALTERNATE INTERIOR angles: 1 below l and right of n,
  % 2 above m and left of n. Each label sits on the bisector of its own sector
  % (l runs at 15.71 deg, n descends at 304.90 deg).
  \node[slab] at ($(P1)+(340.30:0.42)$) {1};
  \node[slab] at ($(P2)+(160.30:0.42)$) {2};
  \node[vlab, right] at (l2) {$l$};
  \node[vlab, right] at (m2) {$m$};
  \node[vlab, below right, xshift=-1pt, yshift=-0pt] at (n2) {$n$};
\end{tikzpicture}
\figcap{1--26}
\end{minipage}\hfill
\begin{minipage}[t]{0.48\linewidth}\centering
\begin{tikzpicture}[scale=1.0]
  \coordinate (l1) at (0,1.35);  \coordinate (l2) at (3.10,1.35);
  \coordinate (n1) at (0.72,0);  \coordinate (n2) at (2.02,2.10);
  \draw[given] (l1) -- (l2);
  \draw[given] (n1) -- (n2);
  \coordinate (Q) at (intersection of l1--l2 and n1--n2);
  \coordinate (P) at ($(n1)!0.30!(n2)$);
  \draw[aux] ($(P)+(-0.66,0)$) -- ($(P)+(1.22,0)$);   % the copied parallel
  \fill (P) circle (1.2pt);
  % as in the book: 2 is the lower-left angle at Q, 1 the upper-right angle at P
  % (alternate interior angles). n rises at 58.24 deg.
  \node[slab] at ($(Q)+(209.12:0.55)$) {2};
  \node[slab] at ($(P)+(29.12:0.55)$) {1};
  \node[vlab] at ($(P)+(299.12:0.34)$) {$P$};
  \node[vlab, left] at (l1) {$l$};
  \node[vlab, above] at (n2) {$n$};
\end{tikzpicture}
\figcap{1--27}
\end{minipage}
\end{bkfigure}}

%--------------------------------------------------------------- Fig 1-28
% Right triangle: right angle at C, hypotenuse c, legs a and b.
\newcommand{\FIGITWENTYEIGHT}{%
\begin{bkfigure}[0.52\linewidth]
\begin{tikzpicture}[scale=1.0]
  \coordinate (A) at (0,0);
  \coordinate (C) at (2.45,0);
  \coordinate (B) at (2.45,1.70);
  \draw[given] (A) -- (C) -- (B) -- cycle;
  \draw[fig] (2.29,0) -- (2.29,0.16) -- (2.45,0.16);
  \node[vlab, left]  at (A) {$A$};
  \node[vlab, below right, xshift=-1pt, yshift=-0pt] at (C) {$C$};
  \node[vlab, above] at (B) {$B$};
  \node[vlab, below] at (1.22,0) {$b$};
  \node[vlab, right] at (2.45,0.85) {$a$};
  \node[vlab, above left, xshift=-0pt, yshift=-1pt] at (1.22,0.85) {$c$};
\end{tikzpicture}
\figcap{1--28}
\end{bkfigure}}


% ---- local macros (hoist into book preamble) ----
% Halftone plates: the book prints full-page photographs facing several pages of
% this chapter. They are reproduced as framed placeholders carrying the original
% caption; see ch01-UNCERTAIN.md.
% Carry the real halftone since 2026-08-17 (Caleb: "I WANT THE ACTUAL HALFTONE
% PICTURES IN THERE"), cut from scan01 pp.19-20 -- see plates/MANIFEST.md.
%   #1 = image file in plates/, #2 = title, #3 = caption
\newcommand{\Iplate}[3]{%
  \par\medskip
  \begin{center}
  \includegraphics[height=64mm]{#1}\par\smallskip
  {\normalfont\textsc{#2}\par\smallskip\footnotesize #3}
  \end{center}
  \par\medskip}
% ---- end local macros ----

\begin{document}
\setcounter{chapter}{1}
\renewcommand{\thefootnote}{\fnsymbol{footnote}}

\chapter*{\raggedright\Huge INTRODUCTION}
\addcontentsline{toc}{chapter}{1.\ Introduction}
\markboth{INTRODUCTION}{}
\vspace{-1em}\noindent\rule{\linewidth}{0.6pt}\vspace{1em}
\figurekey

% The Hilbert plate (leaf facing book p.1) is typeset at the END of the front
% matter (frontmatter.tex, p. xii slot), matching its physical position in the
% book. It was duplicated here until the 2026-08-17 integration pass.

%========================================================= 1-1
\section*{1--1\quad WHY STUDY GEOMETRY?}
\addcontentsline{toc}{section}{1--1\quad Why study geometry?}

\setcounter{footnote}{0}
There are at least four reasons why geometry is studied by so many students in
the high schools of the United States and, in fact, throughout the world. The
first of these reasons is that the logical beauty and precision of geometry has
fascinated men since the time of the ancient Greeks. Geometry is one of the
great achievements of the human mind, and for 2000 years men have regarded its
study as a necessary part of the making of a truly educated person. The feeling
has been that here is at least one thing whose truth is clear, because every
statement can be demonstrated beyond a doubt. People also felt that because the
reasoning in geometry had to be careful and accurate, the study of it would help
one to be careful and accurate in other
activities.\footnote{The Dutch philosopher Benedictus de Spinoza
(1632--1677) was so impressed by the clarity of geometrical reasoning that he
arranged his reasoning on religion in a ``geometric style.''}
Because this great invention of the mind has value as an introduction to correct
logical reasoning, geometry is a required subject for entrance to many colleges
regardless of what one plans to study.

A second important reason for studying geometry is its practical importance.
People in all kinds of occupations have occasional need for a bit of geometry,
and in some fields the study of geometry is a most important step in professional
training. Fields in which it is required are physics, chemistry, engineering,
pure mathematics, statistics, some of the biological sciences, and, more
recently, certain branches of economics and psychology.

Perhaps most students start the study of geometry because of its practical value
or because they are told it is required. But a third reason is that after you
have studied geometry you will have the knowledge that is necessary to give you a
deeper insight into the wonderful details and complexities of our world, both the
natural world and the constructions of mankind. It is this deeper understanding
of the world which can be furthered by the study of geometry that makes its study
important to all people, whether scientists or not. Nature abounds in geometric
designs, and geometric knowledge gives the observer keener interest.

A fourth reason for studying geometry is that we live in what is often called a
scientific age, and even if you yourself are not planning on a career in science
it will be necessary for you to have some understanding and appreciation for the
ways in which a scientist thinks. The study of geometry is a big step toward
getting such understanding.

%========================================================= 1-2
\section*{1--2\quad BRIEF HISTORY OF GEOMETRY}
\addcontentsline{toc}{section}{1--2\quad Brief history of geometry}

\setcounter{footnote}{0}
Knowledge of isolated geometric facts goes back before the beginning of recorded
history. Certainly the early Egyptians and Babylonians (4000--3000 B.C.) knew
many practically useful geometric relations. The construction of the pyramids
alone required considerable practical geometry.

But it was the ancient Greeks who collected the known geometric facts, discovered
new ones, and arranged them into a logically consistent system. (The word
\emph{geometry} comes from two Greek words, \emph{ge} meaning ``earth,'' and
\emph{metrei} meaning ``to measure,'' showing that originally it was thought of
as ``earth measurement.'') This process of organization and discovery took
centuries and occupied the minds of a long list of able men. Remember that the
centuries just before, during, and after the period of greatest Greek political
influence (4th and 5th centuries B.C.) were times of intense intellectual
activity. The leading minds were interested in ideas of all kinds. Mathematics
was but one of their interests.

Some of these early geometers deserve special mention either because of their
mathematical eminence or because of their influence on later men of
philosophy\footnote{The word \emph{philosophy} comes from Greek words meaning
``love of wisdom.'' Look this up in your dictionary. The early Greek
philosophers tried to understand the entire world. More on philosophy will be
found in Chapter 16 of this text.}
or science. One of these early men was Thales of Miletos (6th century B.C.) who
was considered one of the seven wise men of his day. He was the first Greek
mathematician and the first Greek astronomer. Besides knowing some geometric
facts, he probably knew proofs for many of them. The next important Greek for us
was Pythagoras (late 6th century B.C.), who established a school, or brotherhood,
in Croton, Sicily (a part of Greater Greece then). It is often hard to
distinguish whether a discovery originated with Pythagoras or one of his
followers. All ideas were communal property of the brotherhood. Besides the
Pythagorean theorem concerning right triangles, he (or his school) made notable
discoveries in music, astronomy, and arithmetic.

\Iplate{plate-pythagoras}{Pythagoras}{From a fresco by Raphael.\\ Courtesy of \emph{Scripta Mathematica}.}

\Iplate{plate-archimedes}{The Death of Archimedes}{From a mosaic found in the ancient city of
Pompeii.\\ Courtesy of \emph{Scripta Mathematica}.}

\setcounter{footnote}{0}
One Greek geometer in particular deserves mention here because his fame is so
slight compared with his worth. This man is Eudoxus of Cnidus (4th century
B.C.), one of the foremost mathematicians of all time. His fame rests upon his
general theory of proportion, certain geometric constructions, and his ``method
of exhaustion.'' He lived during the time of
Plato.\footnote{Plato and Aristotle are the two most renowned of the Greek
philosophers. Although Eudoxus is less well known, he was a much better
mathematician than either of these men. However, both Plato and Aristotle have
had tremendous influence on later thought. It is said that above the entrance to
Plato's academy were inscribed the words, ``Let no one ignorant of geometry enter
here.''}

We find that the Greeks organized and extended geometry for their own mental
enjoyment. Deductive geometry (that is, with proofs from simple assumptions) is
truly a Greek invention. These men were not seeking practical applications at
all; in fact, an educated Greek regarded applications of geometry as beneath his
dignity. By the time of Plato (4th century B.C.), Greek scholars had a fairly
well-organized body of geometry and perhaps even a geometric text written for
members of Plato's academy.

By far the best known of the ancient Greek geometers is Euclid, who wrote the
\emph{Elements} about 300 B.C. This work is the most successful textbook of all
time and was used throughout the world until well into the present century. It
consists of 13 ``books,'' the first six of which are concerned with plane
geometry. The others deal with arithmetic and solid geometry. The
\emph{Elements} has not only been widely used as written (or as translated); it
has also been the model for innumerable other books. Isaac Newton, the renowned
English mathematician and physicist, wrote his great book, the \emph{Principia},
in the ``geometric style'' even though this often tended to hide the path he took
in making his discoveries. Most of the geometry books used in American high
schools today are somewhat altered and simplified versions of Euclid's
\emph{Elements}. It is astonishing, and a great tribute to Euclid, that this
book has retained its value for over 2000 years.

Little is known of Euclid's life except that he probably learned his mathematics
from pupils of Plato's academy and that he later founded his own school at
Alexandria.\footnote{The city of Alexandria was founded by Alexander the Great in
332 B.C. Alexander was a supporter of learning, as was his successor, Ptolemy I.
It was Ptolemy who started the great library at Alexandria, and it was during his
rule that Euclid opened a school there.}
Undoubtedly he made full use of earlier writings, but the organization of the
\emph{Elements} is his own, as are many significant improvements.

One of the many amusing stories concerning Euclid is the following. Upon being
asked by a student of what profit was the study of geometry, Euclid directed a
slave to ``give the man a coin that he may profit from what he learns.''

\setcounter{footnote}{0}
Many able Greek mathematicians followed Euclid, among whom were Apollonius (3rd
century B.C.) and Archimedes (3rd century B.C.). Archimedes, who had perhaps the
greatest mind of ancient times, did fine work in mathematics, physics, and
engineering. He was employed as a kind of technical consultant in warfare by the
king of Syracuse and was killed by a soldier in the siege of that city in 212
B.C.

With the rise of the Roman Empire, Greece declined as a political power but
retained its eminence as a center of learning and culture well into the Christian
era. Rome added little (compared with Greece) to the arts, science, or
literature. It is in the realm of practical government, warfare, and architecture
that the Romans made their contributions. But when Rome fell to the barbarians
from the north in the 5th century A.D., the ``dark ages'' began. During the dark
ages, mathematical research was largely in the hands of the Arabs. (Our word
\emph{algebra} is of Arabic origin.) Serious scientific work in Europe was almost
at a standstill until the Renaissance, or revival of learning, in the 14th
century.

Beginning with the 16th century, mathematical activity was primarily in fields
other than classical geometry, particularly after the invention of
calculus\footnote{Calculus is a powerful mathematical tool useful in solving
certain kinds of problems. It is a subject required of most science majors in
college.}
by Newton and Leibniz in the 17th century. Nevertheless, significant studies in
geometry were made which culminated in the discovery, in the 19th century by
K.~F. Gauss, N.~Lobachevski, and J.~Bolyai, independently, that geometries other
than Euclid's are possible and just as true (that is, logically correct). It was
during the 19th century that men became aware that there were certain assumptions
made by Euclid which he did not actually state. Therefore Euclidean geometry was
re-examined and made complete and precise according to modern standards. The
first really complete treatment of Euclidean geometry was published by the German
mathematician David Hilbert in 1899. Our book is based upon Hilbert's treatment
(with some simplifications). Thus you will be studying geometry that is almost
2300 years old but dressed up in modern form.

Hilbert (who died in 1943 at the age of 81) had wide-ranging mathematical
interests and made outstanding contributions to many fields. His influence on
American mathematics has been considerable, because at the time when American
mathematicians were beginning to achieve eminence, during the latter part of the
19th century, many of them studied under Hilbert at the University of
G\"ottingen in Germany.

\begin{center}\bfseries Suggested Discussion Topics\end{center}
\nopagebreak
\begin{multicols}{2}
\begin{exlist}
\item Discuss the four reasons listed for studying geometry. Can you think of
      any others?
\item What do you know about Greek civilization other than what has been
      mentioned in this chapter?
\item What were some of the major scientific developments of the 19th century?
      In what countries did they occur?
\item Look up and report on contributions of the Arabs to mathematics.
\item Look up Euclid in an encyclopedia.
\item Look up Pythagoras in an encyclopedia.
\item Look up Archimedes in \emph{The World of Mathematics}, by J.~R. Newman.
\end{exlist}
\end{multicols}

%========================================================= 1-3
\section*{1--3\quad FAMILIAR GEOMETRIC OBJECTS}
\addcontentsline{toc}{section}{1--3\quad Familiar geometric objects}

In your previous study of arithmetic and algebra you have met many ``geometric
objects,'' such as the ones shown in Fig.~1--1, where the examples are all
\emph{plane} figures. Some \emph{solid} figures that you may have encountered are
pictured in Fig.~1--2.

\FIGIONE

\FIGITWO

All these shapes are easily recognized as familiar; you would have no difficulty
in recognizing one of them on sight. However, you may have noticed that we have
not \emph{defined} the shapes, but only drawn them. Part of our job later in this
book will be to give accurate definitions. As you study geometry, it will become
more and more clear that we cannot prove anything without a clear definition. But
in this chapter, we shall most often use drawings to introduce the things you
will study later.

We have used the words ``geometric object'' or ``geometric figure,'' but have
only given examples by drawings on a page of the book. A precise definition is
the following:

\begin{definition}
A geometric figure is any \emph{set of points}.
\end{definition}

Some explanation of the words ``set of points'' is necessary. A set of points is
just a ``bunch'' of points (or ``collection'' of points, or ``aggregate'' of
points). This definition allows some rather strange sets of points to be called
geometric figures, for example those in Fig.~1--3. However, we usually deal with
simpler sets of points than these, namely, with lines, line segments, triangles,
circles, etc. A word of caution is necessary regarding our use of the word
``line.'' By ``line'' we mean a \emph{straight line} indefinite in extent in both
directions. We shall see in Chapter 3 exactly how to state what we need to know
about lines in order to prove other things about them.

\FIGITHREE

\exercisegroup{1--1}
\begin{exercisecols}
\begin{exlist}
\item Describe in your own words what is meant by a \emph{set} of points.
\item Invent some geometric figures of your own and then describe these sets of
      points.
\item Mark a point on a sheet of paper and consider the set of all points on your
      paper such that each point in the set is two inches from this first point.
      Describe this set of points in another way.
\item Some of the shapes drawn in Fig.~1--3 are intended to represent infinitely
      many line segments or circles. Which are they? Can you invent other figures
      with infinitely many segments or circles?
\end{exlist}
\end{exercisecols}

%========================================================= 1-4
\section*{1--4\quad WHAT ARE POINTS?}
\addcontentsline{toc}{section}{1--4\quad What are points?}

Up to now we have been talking about points just as if we knew what they are. But
no one has ever seen a point. The smallest dot you can make on a sheet of paper
is certainly not a point. The dot would look huge to a microbe (if it could see),
and the microbe, which can be seen under a microscope, is very large compared
with the molecules of chemical substances of which it is composed. What, then, is
a point? The answer is that points are an \emph{invention} of the human mind. We
\emph{imagine} that there \emph{ought to be} points, which themselves cannot be
split up into smaller pieces. A given ``hunk'' or portion of space can be split
into two hunks. And each of these can be split in two, and so on indefinitely. A
point is what \emph{should} be left when our hunk has no extension in space or,
to use other words, when it is ``infinitely small.'' Euclid, in fact, gave
something like that as a definition of a point. But such a ``definition'' is not
really a help in geometry because from it we cannot derive anything more. All it
does is make us feel more or less happy about the existence of the things we
\emph{think} we are talking about. The modern point of view (and really Euclid's
too) is that points are undefined. The things that we do define are certain
relations between points or sets of points.

It is essentially the same regarding lines and planes. No one has ever seen a
line or a plane. A thin ray of light or a thin line on the page is only an
approximation to what we imagine that a line should be like. And a very flat
surface, like a mirror, is but an approximation to what we imagine a plane to be.
Therefore right at the beginning of geometry we have to make a strong effort of
the imagination. We must imagine a large set of things, called points, which we
call a plane. Certain parts of this plane, that is, certain special sets of
points in the plane, are called straight lines, or just lines. The study of
geometry begins when we tell \emph{how these points and lines are related to each
other}.

\exercisegroup{1--2}
\begin{exercisecols}
\begin{exlist}
\item Draw two intersecting lines on the blackboard. You will then have some
      chalk hanging on the board. How many reasons can you give to explain that
      what you have drawn are not lines? Does the chalk on both your drawn lines
      correspond to what you think of as a point?
\item Mark two dots on the blackboard. Are they points? Is a ``blob'' of chalk a
      point? How many lines do you feel there should be which lie on two given
      points? If points were blobs of chalk, could you draw more than one line
      through them?
\item When you draw a figure to illustrate a geometric idea, what is the relation
      of the drawn figure to the geometric idea you have in mind?
\item Is it reasonable to expect a drawing to give a proof of something in
      geometry?
\item If an object had length or width, would you call it a point? Can you see
      points and lines?
\item Can you describe what you mean by straightness of a line?
\end{exlist}
\end{exercisecols}

%========================================================= 1-5
\section*{1--5\quad TOOLS YOU WILL NEED}
\addcontentsline{toc}{section}{1--5\quad Tools you will need}

\setcounter{footnote}{0}
Even though points and lines are creations of the mind it is necessary for us to
draw pictures of them in order to see what is ``going on.'' Thus you will need a
\emph{ruler} to draw pictures of straight lines. A ruler without a marked scale
on it is called a \emph{straightedge}, and for the constructions we shall make
this is all that is absolutely necessary. But it is very convenient to be able to
\emph{measure} the length of a line segment, and for this reason an ordinary foot
ruler, or a six-inch one, will do. However, we shall not actually deal with
measurement\footnote{It is interesting and surprising that Euclid in his
\emph{Elements} did not discuss measurement, that is, the process of getting a
\emph{number} for the distance between two points.}
of lengths until Chapter 4.

Angles will be discussed more fully later, and we shall consider measurement of
angles in Chapter 12. But you are already familiar with a device, called the
\emph{protractor}, for measuring pictures of angles drawn on paper or on the
blackboard.

You will also need a device, called a \emph{compass}, for drawing pictures of
circles. Undoubtedly you are familiar with its use.

With these simple tools you can make a number of ``constructions.'' These
constructions will seem to be obviously correct. However, we will prove later in
the book that they actually are correct.

You have learned in your previous mathematics how to use these tools. The
exercises below are for review of definitions and practice in the use of ruler,
compass, and protractor. However, throughout this course your best tool will be
your ability to think straight.

\exercisegroup{1--3}
% multicols dropped for this group: it carries figures, and the book sets each
% one immediately after the exercise that cites it, at the full measure.
\begin{exercisecols}
\begin{exlist}
\item Draw a picture of a line, using your straightedge. Now if you sharpened
      your pencil \emph{very} sharp, could you draw a picture of a line inside the
      first one? Was your drawing really what you intended a line to be?
\item Draw a picture of a circle with your compass. How many differences can you
      find between your drawing and an imaginary perfect circle?
\item Draw a picture of an angle. Measure it with your protractor. Have a friend
      measure the angle. Do you get \emph{exactly} the same answers? How accurate
      do you think your measurement is?
\item How long are the sides of an angle? When you draw a picture of an angle,
      you cannot draw the sides infinitely long because your straightedge and
      paper are limited in size.
\item Draw a picture of a ray from a point $A$. Draw another ray from $A$. What
      we call the \emph{angle} consists of the points on the two rays. The points
      between the rays form the \emph{interior} of the angle. The other points of
      the plane form the \emph{exterior} of the angle. (Fig.~1--4.)
\exfig{\FIGIFOUR}
\item A \emph{right angle} has a measure of $90^\circ$. Construct a right angle
      with your protractor.
\item Draw an angle of $37^\circ$, of $120^\circ$, of $85^\circ$.
\item An angle is \emph{bisected} if a ray from the vertex splits the angle into
      two angles of the same size. In Fig.~1--5, the ray $AC$ bisects
      $\ang{BAD}$ because $\ang{CAB}$ and $\ang{DAC}$ are the same size. The ray
      $AC$ is the bisector of the angle. Draw an angle of $45^\circ$ and bisect
      it, using your protractor.
\exfig{\FIGIFIVE}
\item Draw a picture of a triangle. Measure each of the angles of the triangle.
      What is their sum? Repeat for several different triangles. Can you guess
      something that is true for all triangles?
\item As in Fig.~1--6, an \emph{acute angle} is an angle of less than
      $90^\circ$; an \emph{obtuse angle} is an angle of more than $90^\circ$.
      Draw two acute angles and two obtuse angles with your protractor.
\exfig{\FIGISIX}
\item Draw a triangle and measure the sides and the angles. Which side is
      opposite the largest angle? opposite the smallest angle? Draw another
      triangle and repeat the measurements. Can you guess something which now
      seems to be true for all triangles?
\item A triangle that has one right angle is called a \emph{right triangle}. Draw
      a picture of a right triangle, as in Fig.~1--7. Measure the two acute
      angles of the triangle. What is their sum? Draw another right triangle and
      make the same measurements. Is the sum of the two angles the same as
      before?
\exfig{\FIGISEVEN}
\item One angle is said to be \emph{complementary} to another if their sum is
      $90^\circ$. Construct two such angles with your protractor. In Fig.~1--8,
      if $\ang{AOB}$ is a right angle, what can you say about $\ang{COA}$? about
      $\ang{COB}$?
\exfig{\FIGIEIGHT}
\item One angle is said to be \emph{supplementary} to another if their sum is
      $180^\circ$. Construct two such angles with your protractor. What can you
      say about $\ang{AOB}$ and about $\ang{BOC}$ in Fig.~1--9?
\exfig{\FIGININE}
\item Two distinct rays $OA$ and $OB$ which lie on one straight line form a
      \emph{straight angle} (Fig.~1--10). What is the measure of a straight
      angle?
\exfig{\FIGITEN}
\item A \emph{quadrilateral} (Fig.~1--11) is a figure with four sides, just as a
      triangle is a figure with three sides. Draw a picture of a quadrilateral and
      measure its angles. What is their sum?
\exfig{\FIGIELEVEN}
\item Add the following:
  \begin{enumerate}[label=(\alph*), leftmargin=1.9em, itemsep=1pt, topsep=2pt]
  \item $12^\circ\ 10'\ 25''$, $31^\circ\ 48'\ 30''$, and $43^\circ\ 50'\ 40''$.
  \item $17^\circ\ 18'\ 19''$, $90^\circ\ 9'\ 9''$, and $41^\circ\ 48'\ 48''$.
  \item $121^\circ\ 18'\ 17''$, $40^\circ\ 40'\ 40''$, and $18^\circ\ 1'\ 3''$.
  \item $66^\circ\ 55'\ 44''$, $55^\circ\ 44'\ 33''$, and $44^\circ\ 33'\ 22''$.
  \end{enumerate}
\item Subtract:
  \begin{enumerate}[label=(\alph*), leftmargin=1.9em, itemsep=1pt, topsep=2pt]
  \item $41^\circ\ 15'\ 13''$ from $63^\circ\ 28'\ 44''$.
  \item $41^\circ\ 15'\ 13''$ from $50^\circ\ 18'\ 7''$.
  \item $61^\circ\ 28'\ 37''$ from $80^\circ\ 18'\ 23''$.
  \end{enumerate}
\item What is the complement of an angle of
  \begin{enumerate}[label=(\alph*), leftmargin=1.9em, itemsep=1pt, topsep=2pt]
  \item $51^\circ\ 17'\ 18''$?
  \item $42^\circ\ 38'\ 29''$?
  \item $A^\circ$?
  \end{enumerate}
\item What is the supplement of an angle of
  \begin{enumerate}[label=(\alph*), leftmargin=1.9em, itemsep=1pt, topsep=2pt]
  \item $51^\circ\ 17'\ 18''$?
  \item $138^\circ\ 4'\ 5''$?
  \item $A^\circ$?
  \end{enumerate}
\item If an angle is twice as large as its complement, what is its degree
      measure? [\emph{Hint}: Set up an equation to solve.]
\item Two angles have a sum of $56^\circ\ 28'$ and a difference of
      $8^\circ\ 46'$. Find the angles.
\item Find the complement and supplement of an angle of $56^\circ\ 28'\ 31''$.
      What is the difference between the supplement and complement? Can you guess
      something that is true for all acute angles?
\item What kind of angle is the supplement of an acute angle? of an obtuse angle?
      Does an obtuse angle have a complement?
\end{exlist}
\end{exercisecols}

%========================================================= 1-6
\section*{1--6\quad RULER AND COMPASS}
\addcontentsline{toc}{section}{1--6\quad Ruler and compass}

During their study of geometry, the ancient Greeks found that many interesting
figures could be drawn with a straightedge and compass alone. Drawings needing
only these tools are called \emph{ruler and compass constructions}. The Greeks
tried to solve all construction problems using only these tools, but they
encountered some that could not be done this way.

The exercises below involve some simple ruler and compass constructions that the
Greeks knew. Exercise 2, in fact, is the first proposition in Euclid's
\emph{Elements}. We shall return to these constructions again in Chapter 7, where
we shall \emph{prove} that some of them are correct.

\exercisegroup{1--4}
% multicols dropped for this group: it carries figures, and the book sets each
% one immediately after the exercise that cites it, at the full measure.
\begin{exercisecols}
\begin{exlist}
\item A triangle having two sides of equal length is called an \emph{isosceles
      triangle}. Draw a picture of an isosceles triangle. Measure the angles
      opposite the equal sides. Draw several isosceles triangles, each with the
      equal sides three inches long.
\item A triangle with all sides the same length is called an \emph{equilateral
      triangle} (Fig.~1--12). Draw a picture of an equilateral triangle. This is
      easy to do with your compass and straightedge. What is the degree measure of
      each angle in an equilateral triangle? A triangle with all angles equal is
      called an \emph{equiangular triangle}. Are equilateral triangles
      equiangular? Are equiangular triangles equilateral?
\exfig{\FIGITWELVE}
\item Draw an angle like that in Fig.~1--13. Use your compass to copy $\ang{AOB}$
      as $\ang{A'O'B'}$ in the figure. Measure the two angles with your protractor
      to see if they are the same size.
\exfig{\FIGITHIRTEEN}
\item Draw the figures shown in Fig.~1--14.
\exfig{\FIGIFOURTEEN}
\item Use your straightedge and compass to draw and bisect an angle
      (Fig.~1--15). Check your work by measurement with a protractor. State in
      your own words why you think the construction is correct.
\exfig{\FIGIFIFTEEN}
\item Angles which have a common vertex and a common side and have no interior
      points in common are said to be \emph{adjacent}. Thus, in Fig.~1--16,
      $\ang{AOB}$ is adjacent to $\ang{BOC}$. Is $\ang{AOC}$ adjacent to
      $\ang{BOC}$? Are the angles in Fig.~1--17 adjacent to one another?
\exfig{\FIGISIXTEEN}
\exfig{\FIGISEVENTEEN}
\item Draw a picture of a segment. From some point $A$ on a line, construct with
      your compass a segment $AB$ as long as the segment you first drew. (See
      Fig.~1--18.)
\exfig{\FIGIEIGHTEEN}
\item Two lines are \emph{perpendicular} if they intersect at right angles.
      Another way of defining perpendicular would be to say that all the adjacent
      angles formed by the intersecting lines are the same size. To indicate that
      $l$ and $m$ are perpendicular, we write $l \perp m$.

      A line $l$ is called the \emph{perpendicular bisector} of the segment $AB$
      if the line $l$ is perpendicular to $AB$ and intersects the segment $AB$ in
      a point $C$ such that $AC$ and $CB$ have equal lengths. The point $C$ is
      called the \emph{midpoint} of $AB$. (See Fig.~1--19.)

      Draw a line segment and construct its perpendicular bisector with
      straightedge and compass. Check by measurement. State why you think the
      construction is correct.
\exfig{\FIGININETEEN}
\item Draw a line $l$ and mark a point $O$ on it. With straightedge and compass,
      construct a line $m \perp l$ and passing through $O$. (See Fig.~1--20.)
\exfig{\FIGITWENTY}
\item Draw a line $l$ and mark a point $P$ not on the line. With your straightedge
      and compass, construct a line $m$ through $P$ perpendicular to $l$. Check by
      measurement with your protractor. State why you think the construction is
      correct. (See Fig.~1--21.)
\exfig{\FIGITWENTYONE}
\item A line segment drawn from a vertex of a triangle to the midpoint of the
      opposite side is called a \emph{median} of the triangle. A triangle has
      three medians that meet in one point $O$ (see Fig.~1--22). We shall prove
      this later.

      Draw a triangle and its medians. Can you describe where the medians
      intersect? Try to find how to describe $O$ by measurement. Make this
      construction for several triangles.
\exfig{\FIGITWENTYTWO}
\item Cut a triangle out of cardboard. Find where the medians intersect. If your
      work is accurate you will find that the triangle will balance on your finger
      at the place where the medians intersect (Fig.~1--23).
\exfig{\FIGITWENTYTHREE}
\item An \emph{altitude} of a triangle is a line segment drawn from a vertex
      perpendicular to the line \emph{containing} the opposite side [see
      Fig.~1--24, (a) and (b)]. A triangle has three altitudes. The lines
      containing the altitudes of a triangle always intersect in one point [see
      Fig.~1--24(c)]. Draw several triangles and their altitudes.
\exfig{\FIGITWENTYFOUR}
\item Draw a triangle and bisect each of the angles of the triangle. The three
      bisectors will meet in one point (Fig.~1--25).
\exfig{\FIGITWENTYFIVE}
\item Two lines, $l$ and $m$, are said to be \emph{parallel} if they do not
      intersect. We indicate ``$l$ is parallel to $m$'' by writing $l \pll m$.
      (Remember that lines are infinitely long.) If another line $n$ intersects
      the parallel lines, then the angles 1 and 2 in Fig.~1--26 are the same size.
      Conversely, if angles 1 and 2 are the same size, then the lines are
      parallel.

      As in Fig.~1--27, draw a line $l$ and a line $n$ intersecting it. Choose a
      point $P$ on $n$ but not on $l$. At $P$, copy angle 2 in order to get a line
      parallel to $l$. Make this construction with your protractor and then repeat
      it with your straightedge and compass.

      We shall prove later that this construction with straightedge and compass is
      correct.
\exfig{\FIGITWENTYSIX}
\exfig{\FIGITWENTYSEVEN}
\item Draw a line $l$ and mark a point $A$ not on $l$. Construct a line through
      the point $A$ parallel to $l$.
\item Draw a right triangle (Fig.~1--28). The side opposite the right angle is
      called the \emph{hypotenuse}. The other two sides are called the \emph{legs}
      of the triangle. Measure the lengths of the sides in your right triangle. If
      $a$, $b$, and $c$ are the lengths of the sides, where $c$ is the hypotenuse,
      you will find that $a^2 + b^2 = c^2$. This formula is called the
      \emph{Pythagorean theorem}. It was discovered (probably) by the Greek
      Pythagoras in the 6th century B.C. We will prove this important theorem
      later in the book. Test the truth of the Pythagorean theorem by measuring
      several different right triangles.
\exfig{\FIGITWENTYEIGHT}
\end{exlist}
\end{exercisecols}

%========================================================= Review
\section*{REVIEW OF CHAPTER 1}
\addcontentsline{toc}{section}{Review of Chapter 1}
\markboth{REVIEW OF CHAPTER 1}{}

\begin{exlist}
\item Make sure that you are familiar with all the following terms:
\end{exlist}

\begin{center}
\begin{tabular}{@{}p{0.42\linewidth}p{0.42\linewidth}@{}}
line segment            & perpendicular    \\
angle                   & right angle      \\
angle bisector          & equilateral      \\
triangle                & equiangular      \\
degree, minute, second  & isosceles        \\
midpoint                & median           \\
complementary           & altitude         \\
supplementary           & hypotenuse       \\
straight angle          & diameter         \\
obtuse angle            & radius           \\
acute angle             & adjacent angles  \\
ray                     & intersect        \\
set                     & contain          \\
\end{tabular}
\end{center}

\begin{exlist}\setcounter{exlisti}{1}
\item Make sure that you can do the following straightedge and compass
      constructions:
  \begin{enumerate}[label=(\alph*), leftmargin=2.2em, itemsep=1pt, topsep=2pt]
  \item Copy a segment.
  \item Draw the bisector of a segment.
  \item Draw a line $\perp$ to a given line from a point not on the given line.
  \item Draw a line $\perp$ to a given line from a point on a given line.
  \item Draw a line $\pll$ to a given line, through a point not on the line.
  \item Bisect an angle.
  \item Copy an angle.
  \end{enumerate}
\end{exlist}

\begin{center}\bfseries Brief Review Tests\end{center}
\nopagebreak

Decide whether the following statements are true or false. If in doubt, look up
the answer. Do not guess.

\medskip
\noindent\textbf{I.} \emph{History}
\begin{exlist}
\item Euclid lived before Pythagoras.
\item Euclid collected, organized, and added to the geometry of his day.
\item Euclid's geometry was incomplete as he wrote it. It was completed in the
      19th century.
\item Archimedes lived at the same time as Euclid.
\item Thales lived before Euclid.
\item Greece was superior to Rome in all ways.
\item Rome was a great center for the study of geometry.
\end{exlist}

\medskip
\noindent\textbf{II.} \emph{Geometry}
\begin{exlist}
\item A straightedge is used for measurement.
\item A line may have finite length.
\item A straight angle contains $90^\circ$.
\item If angles have a common vertex and a common side, they are adjacent.
\item Complementary angles have a sum of $180^\circ$.
\item In a right triangle the side opposite the right angle is called the
      hippopotamus.
\item The altitude of a triangle is its distance from the center of the earth.
\item If two lines intersect they are parallel.
\end{exlist}

\begin{center}\bfseries ARITHMETIC AND ALGEBRA REFRESHER\end{center}
\nopagebreak

\begin{exlist}
\item Carry out the arithmetic operations indicated:
  \begin{multicols}{2}
  \begin{enumerate}[label=(\alph*), leftmargin=2.2em, itemsep=1pt, topsep=2pt]
  \item $[(12)(8)] \div 6$
  \item $2992 \div 68$
  \item $326 + (-257)$
  \item $\frac{1}{2} + \frac{1}{4}$
  \item $\frac{1}{3} \div \frac{11}{3}$
  \item $56789 \div 357$
  \item $\frac{2}{3} - \frac{2}{5}$
  \item $3.762 \div 5.7$
  \item $\frac{1}{2} \cdot \frac{1}{3}$
  \item $\frac{1}{2} \div \frac{1}{3}$
  \item $(-24) \div (-6/5)$
  \item $[(-6.1) + (-8.5)] - 2.7$
  \item $3/7 \cdot 35/24$
  \item $2.5679/0.1372$
  \item $[3.4 - 8.5] + (-2.7)$
  \item $(5 + \frac{1}{2})/33$
  \end{enumerate}
  \end{multicols}
\item Write the answers to the following. If you take more than a minute and miss
      any part of this problem, you need more review.
  \begin{multicols}{3}
  \begin{enumerate}[label=(\alph*), leftmargin=2.2em, itemsep=1pt, topsep=2pt]
  \item $2 \times 2 =$
  \item $2 \div 2 =$
  \item $2 \times 0.2 =$
  \item $2 \div 0.2 =$
  \item $2 \div \frac{1}{2} =$
  \item $2 \times \frac{1}{2} =$
  \item $\frac{1}{2} + \frac{1}{3} =$
  \end{enumerate}
  \end{multicols}
\item Write formulas for the following. Be sure that you tell what each letter
      stands for. For example: If $A$ is the area of a rectangle $l$ its length
      and $w$ its width, then $A = lw$.
  \begin{enumerate}[label=(\alph*), leftmargin=2.2em, itemsep=1pt, topsep=2pt]
  \item The volume of a box.
  \item The area of a circle.
  \item The circumference of a circle.
  \item The area of a triangle.
  \item The product of a number $z$, by one more than the number.
  \item Twice a number plus 5 is equal to 13.
  \item One number is 4 times another.
  \end{enumerate}
\item Find $x$ if
  \begin{multicols}{2}
  \begin{enumerate}[label=(\alph*), leftmargin=2.2em, itemsep=1pt, topsep=2pt]
  \item $2x - 5 = 3$
  \item $3x - 5 = x + 7$
  \item $2x = 3$
  \item $ax = b$, $a \neq 0$
  \item $x^2 = 4$
  \item $\frac{1}{3} \cdot (12x) = 44$
  \item $2x/3 + 1 = 5$
  \item $4x + 5x - x = 96$
  \item $x \div 6 = 5 \div 24$
  \item $6.4 = 2x - 0.4$
  \end{enumerate}
  \end{multicols}
\item If $x$ and $y$ are numbers, what do we call $xy$? $x + y$?
\item Write equations for the following statements:
  \begin{enumerate}[label=(\alph*), leftmargin=2.2em, itemsep=1pt, topsep=2pt]
  \item The sum of two numbers is 17 and their difference is 7.
  \item A number $z$ is equal to the sum of the product of $x$ and $y$, and the
        product of $a$ and $b$.
  \item The number $\pi$ is equal to the quotient of the circumference of a
        circle divided by its diameter.
  \item The volume of a sphere is equal to $\frac{4}{3}$ of the product of $\pi$
        and the cube of the radius.
  \item Einstein's famous formula for the equivalence of mass and energy is as
        follows: If $E$ is the energy of a mass $m$, and $c$ is the velocity of
        light, in suitable units, then $E$ is equal to the product of the mass by
        the square of the velocity of light. Write this formula.
  \end{enumerate}
\end{exlist}

% The Aristotle plate (book p.20) is NOT emitted here. It is the facing-page
% plate for the Chapter 2 opener and ch02.tex already emits it immediately
% before \chapter*{LOGIC}. Keeping it in both units printed it twice in the
% assembled book. Ownership: ch02. See ch01-UNCERTAIN.md "Text review".

\end{document}
